% =========================
% Problem Set Template (LaTeX)
% =========================
\documentclass[11pt]{article}

% ---------- Page + layout ----------
\usepackage[margin=1in]{geometry}
\usepackage{setspace}
\setstretch{1.1}
\usepackage{parskip} % space between paragraphs, no indent

% ---------- Math ----------
\usepackage{amsmath, amssymb, amsfonts, amsthm, mathtools}

% ---------- Lists + links ----------
\usepackage{enumitem}
\usepackage[hidelinks]{hyperref}

% ---------- Header / footer ----------
\usepackage{fancyhdr}
\pagestyle{fancy}
\fancyhf{}
\lhead{\textbf{\course}}
\chead{\textbf{\pset}}
\rhead{\textbf{\name}}
\cfoot{\thepage}

% ---------- Optional: nicer tables ----------
\usepackage{booktabs}

% =========================
% Metadata (edit these)
% =========================
\newcommand{\name}{Tate Mason}
\newcommand{\course}{Econometrics II}
\newcommand{\pset}{Problem Set 1}
\newcommand{\duedate}{February 6, 2026}

% =========================
% Question / solution macros
% =========================
\newcounter{q}
\renewcommand{\theq}{\arabic{q}}

% Usage: \question{Title}
\newcommand{\question}[1]{%
  \refstepcounter{q}
  \vspace{0.75em}
  \noindent{\Large\textbf{Question \theq.} \normalsize\textbf{#1}}%
  \vspace{0.25em}
  \par
}

% Usage: \subquestion{Title}  -> (a), (b), ...
\newcounter{subq}[q]
\renewcommand{\thesubq}{\alph{subq}}
\newcommand{\subquestion}[1]{%
  \refstepcounter{subq}
  \vspace{0.5em}
  \noindent{\textbf{(\thesubq)} \textbf{#1}}%
  \par
}

% Solution environment: comment out \begin{solution}...\end{solution} if not needed
\newenvironment{solution}{%
  \vspace{0.25em}
  \noindent\textit{Solution. }\ignorespaces
}{\vspace{0.5em}\par}

% =========================
% Easy math macros
% =========================
% Common sets
\newcommand{\R}{\mathbb{R}}
\newcommand{\N}{\mathbb{N}}

% Operators
\DeclareMathOperator{\E}{\mathbb{E}}
\DeclareMathOperator{\Var}{Var}
\DeclareMathOperator{\Cov}{Cov}
\DeclareMathOperator{\plim}{plim}
\DeclareMathOperator*{\argmin}{argmin}
\DeclareMathOperator*{\argmax}{argmax}

% Probability
\newcommand{\P}{\mathbb{P}}

% Vectors / matrices (bold)
\newcommand{\vct}[1]{\mathbf{#1}}
\newcommand{\mtx}[1]{\mathbf{#1}}

% Derivatives
\newcommand{\dd}{\,\mathrm{d}}
\newcommand{\pd}[2]{\frac{\partial #1}{\partial #2}}
\newcommand{\pdd}[2]{\frac{\partial^2 #1}{\partial #2^2}}
\newcommand{\od}[2]{\frac{\mathrm{d} #1}{\mathrm{d} #2}}
\newcommand{\odd}[2]{\frac{\mathrm{d}^2 #1}{\mathrm{d} #2^2}}

% Norms / inner products
\newcommand{\norm}[1]{\left\lVert #1 \right\rVert}
\newcommand{\ip}[2]{\left\langle #1, #2 \right\rangle}

% Convenience
\newcommand{\1}{\mathbf{1}} % indicator (use \1\{A\} with braces)
\newcommand{\eps}{\varepsilon}

% =========================
% Document
% =========================
\begin{document}

% ---------- Title block ----------
{\Large\textbf{\course} \hfill \textbf{\pset}}\\
\textbf{Name:} \name \hfill \textbf{Due:} \duedate\\
\rule{\textwidth}{0.4pt}

\question{Example Code and Walkthrough}
  \subquestion{- (e): Example and Walkthrough}
  \subquestion{Is it surprising that the heteroskedastic and homoskedastic starndart errors for the OLS are so similar?}
  \begin{solution}
  \end{solution}
  \subquestion{What is the additional parameter we are estimating in part (d) using the likelihood based approach?}
  \begin{solution}
  \end{solution}
  \subquestion{Show analytically why the MLE and MoM approaches yield identical parameter estimates to OLS}
  \begin{solution}
  \end{solution}
  \subquestion{What computational method is "fminunc" using to find the minimum of the function (assuming the analytical gradient is not supplied)?}
  \begin{solution}
  \end{solution}

\question{Solo Exercise}
  \subquestion{Estimate $\beta$ and its s.e. using MLE with the available data}
    \begin{solution}

      \begin{centering}
        \begin{tabular}{ll}
          Parameter & Estimate \\
          \hline

          $\hat{\beta}_1$ & 0.519 \\
                  & (0.016) \\
          $\hat{\beta}_2$ & 0.258 \\
                    & (0.013) \\
          $\hat{\beta}_3$ & 0.485 \\
                    & (0.008) \\
          $\hat{\beta}_4$ & 0.343 \\
                    & (0.005) \\
        \end{tabular}
      \end{centering}

    \end{solution}
  \subquestion{Estimate $\beta$ using NLS. Can the inverse of the numerical Hessian from "fminunc" be used to construct s.e.? (Hint: 417 of Wooldridge)}
  \subquestion{Calculate s.e. for NLS using the analytical formulas in eq. 12.52 of Wooldridge (desc. below said equation).}
  \subquestion{Calculate s.e. for NLS using eq. 12.48 of Wooldridge. Calculate the score analytically, but use the Hessian provided by "fminunc". Why can't the gradient from "fminunc" be used?}
  \subquestion{Calculate s.e. for NLS using bootstrap (iter = 500). Note the command datasample() command can be helpful.}
  \subquestion{Why are s.e. for MLE significantly smaller than NLS?}
  \subquestion{If analytical score is too hard to construct, and bootstrap too time consuming, what can be done to construct s.e.?}
\end{document}
